% ID: SEL-D0033-OB-E03
% ESTADO: cotejada. Segunda pasada visual y borrador compilado revisados.
% Fuente: Selectividad/matematicas_aplicadas/titular_septiembre/Examen-Matemáticas Aplicadas a las Ciencias Sociales.pdf
% SHA-256: 0f95b6aefa78f1a13e2115b281338bd807177671c303c380180957fe03f0a605
% Página: 2; opción B; ejercicio original 3.

Un hotel dispone de tres lavadoras industriales $L_1$, $L_2$ y $L_3$ para el servicio de lavandería. El 50\% de los lavados los realiza $L_1$, el 30\% los hace $L_2$ y el resto $L_3$. La lavadora $L_1$ produce un 5\% de lavados defectuosos, $L_2$ produce un 15\% y $L_3$ un 20\%. Se elige al azar un lavado del hotel.
\begin{enumerate}[label=\alph*)]
\item \textbf{(1.5 puntos)} Calcule la probabilidad de que no sea defectuoso.
\item \textbf{(1 punto)} Calcule la probabilidad de que el lavado haya sido realizado por $L_1$, sabiendo que ha sido defectuoso.
\end{enumerate}
